\documentclass[12pt]{book}
\newcommand{\hmmax}{0}
\newcommand{\bmmax}{0}

\usepackage[small,compact]{titlesec}
\usepackage{titletoc,titleps}

\usepackage{kpfonts}
\usepackage{bm,yfonts}
\usepackage[bb=boondox,bbscaled=1.05]{mathalfa}
%\usepackage{BOONDOX-ds}
%\usepackage{MnSymbol}
%\usepackage{unicode-math}





\usepackage{mathtools}
\usepackage[standard,thmmarks,amsmath]{ntheorem}
\usepackage{cleveref}
%\usepackage[capitalize]{cleveref}
\usepackage{stackrel}
\usepackage{enumitem}
\usepackage{graphicx}
\usepackage{fancyhdr}
\usepackage{color}
\usepackage[font=small]{caption}
\usepackage[margin=3cm]{geometry}
% \makeatletter
%   \@addtoreset{equation}{\thechapter}
%   \@addtoreset{proposition}{\thechapter}
% \makeatother  
   
\usepackage[numbers]{natbib}
\bibliographystyle{./plainnat_modified}

\usepackage{tikz}
\usetikzlibrary{positioning}
\usetikzlibrary{decorations.pathreplacing,calc} 
\usetikzlibrary{shapes.geometric}
%\usepackage{palatino}
 
%\pagestyle{fancy}
%\fancyhf{}
%\lhead{553.740: Introduction to Machine Learning}
%\rhead{Page \thepage}
%\cfoot{{\color{red} FOR CLASS USE ONLY. DO NOT DISTRIBUTE.}}


%\usepackage[tt]{titlesec}


%\renewcommand{\mathbb}{\mathbbm}
%\titleformat{\subsubsection}[runin]
\input Headers/newcom
\input Headers/ml_preamble
\input Headers/question_preamble



%\newcommand{\indp}{{\perp\nolinebreak\!\!\!\!\perp}}
\newcommand{\indp}{{\Bot}}
\newcommand{\cind}[3]{(#1\indp#2\mid #3)}
\newcommand{\twop}[1]{{\boldsymbol{\mathcal P}({#1})}}
\newcommand{\osc}{\mathit{osc}}
%\renewcommand{\hat}{\widehat}
%\renewcommand{\bar}{\overline}
%\renewcommand{\tilde}{\widetilde}
\solutionsfalse
\uselabeltrue

\newcommand{\Eq}[1]
{
\[
\displaystyle #1
\]
}

\newcommand{\alert}[1]{{\em #1}}

\theoremstyle{plain}
\theorembodyfont{\em}
%\renewtheorem{definition}[chapter]{Definition}
%\renewtheorem{theorem}[chapter]{Theorem}
%\newtheorem{result}{Result}[chapter]
%\renewtheorem{proposition}[chapter]{Proposition}
%\newtheorem{conditions}{Conditions}[chapter]
%\newtheorem{hypothesis}{Hypothesis}[chapter]
%\renewtheorem{corollary}{Corollary}[chapter]
%\renewtheorem{lemma}[chapter]{Lemma}
%\newtheorem{properties}{Properties}[chapter]
%\newtheorem{property}{Property}[chapter]
\newtheorem{all}{all}[chapter]
\renewtheorem{definition}[all]{Definition}%[chapter]
\renewtheorem{theorem}[all]{Theorem}%[chapter]
\newtheorem{result}[all]{Result}%[chapter]
\renewtheorem{proposition}[all]{Proposition}%[chapter]
\newtheorem{conditions}[all]{Conditions}%[chapter]
\newtheorem{hypothesis}[all]{Hypothesis}%[chapter]
\renewtheorem{corollary}[all]{Corollary}%[chapter]
\renewtheorem{lemma}[all]{Lemma}%[chapter]
\newtheorem{properties}[all]{Properties}%[chapter]
\newtheorem{property}[all]{Property}%[chapter]

\theorembodyfont{\rm}
\theoremsymbol{\ensuremath{\blacklozenge}}
\renewtheorem{remark}[all]{Remark}%[chapter]
\newtheorem{notation}[all]{Notation}%[chapter]
\renewtheorem{example}[all]{Example}%[chapter]


\theoremstyle{break}
\theoremsymbol{}
\theoremprework{\bigskip\hrule}
\theorempostwork{\hrule\bigskip}
\newtheorem{algorithm}{Algorithm}[chapter]

\crefname{equation}{}{}
\Crefname{equation}{Equation}{Equations}
\crefname{algorithm}{Algorithm}{Algorithms}
\Crefname{algorithm}{Algorithm}{Algorithms}
\setlist[enumerate]{label=\arabic*., wide=0.5cm}

%\providecommand{\citet}{\cite}
%\providecommand{\citep}{\cite}
\renewcommand{\cite}{\citep}
\graphicspath{{./},{Figures/}}

\setlist[enumerate]{label= (\arabic*), labelsep=*, leftmargin=0.25cm, wide=0cm, itemsep=0cm, topsep=0.1cm, ref=\arabic*}
\setlist[itemize]{label= \textbullet, labelsep=*, leftmargin=0.25cm, wide=0cm, itemsep=0cm, topsep=0.1cm}

\newcommand{\problems}[1]{\newpage
\markright{PROBLEMS}
\section*{Problems}
\addcontentsline{toc}{section}{Problems}
\input #1
}

% comment out to include problems
\renewcommand{\problems}[1]{}

%\newcommand{\pe}[2]{#1^{(#2)}}

%\numberwithin{section}{chapter}
%\numberwithin{subsection}{chapter}

\title{Introduction to Machine Learning}
\author{Laurent Younes}

\begin{document}

\maketitle
\tableofcontents

\setlength{\parskip}{1em}

\input preface
\input notation
\input linalg

\input optimization
%\newpage
%\markright{PROBLEMS}
%\section*{Problems}
%\addcontentsline{toc}{section}{Problems}
\problems{Problems_Optimization}


\input intro
\input errors
\problems{Problems_Bayes_Rule}


\input higher
%\newpage
%\markright{PROBLEMS}
%\section*{Problems}
%\addcontentsline{toc}{section}{Problems}
\problems{Problems_Kernels}


\input linear_regression
%\newpage
%\markright{PROBLEMS}
%\section*{Problems}
%\addcontentsline{toc}{section}{Problems}
\problems{Problems_Linear_Regression}

\input linear_classification
%\newpage
%\markright{PROBLEMS}
%\section*{Problems}
%\addcontentsline{toc}{section}{Problems}
\problems{Problems_Classification}

\input nn
%\newpage
%\markright{PROBLEMS}
%\section*{Problems}

\input tree
%\newpage
%\markright{PROBLEMS}
%\section*{Problems}

\input neuralNets
%\newpage
%\markright{PROBLEMS}
%\section*{Problems}
%\addcontentsline{toc}{section}{Problems}
\problems{Problems_Neural_Networks}

\input mcmc
\input sgd_markov

\input mrf

\input parametric_inference

\input bayesian_networks


\input variationalBayes

\input learning_graphical

%\newpage
%\markright{PROBLEMS}
%\section*{Problems}

\input generative

\input clustering
%\newpage
%\markright{PROBLEMS}
%\section*{Problems}
%\addcontentsline{toc}{section}{Problems}
\problems{Problems_Clustering}

\input dim-red
%\newpage
%\markright{PROBLEMS}
%\section*{Problems}
%\addcontentsline{toc}{section}{Problems}
\problems{Problems_Factor_Analysis}

\input manifold
%\newpage
%\markright{PROBLEMS}
%\section*{Problems}
%\addcontentsline{toc}{section}{Problems}
\problems{Problems_Manifold_learning}

\input ms
%\newpage
%\markright{PROBLEMS}
%\section*{Problems}
%\addcontentsline{toc}{section}{Problems}
\problems{Problems_Concentration}

\bibliography{machine_learning}

\end{document}

\chapter{Problems with Solutions}
\setcounter{chapter}{1}
\solutionstrue
\uselabelfalse

\chapter{Problems with solutions}
\markright{PROBLEMS WITH SOLUTIONS}
\section*{Basic concepts}
\stepcounter{chapter}
\input Problems_Bayes_Rule
\section*{Matrix Analysis}
\stepcounter{chapter}
\section*{Reproducing Kernels}
\stepcounter{chapter}
\input Problems_Kernels
\section*{Optimization}
\stepcounter{chapter}
\input Problems_Optimization
\section*{Linear Regression}
\stepcounter{chapter}
\input Problems_Linear_Regression
\section*{Linear Classification}
\stepcounter{chapter}
\input Problems_Classification
\section*{Nearest Neighbors}
\stepcounter{chapter}
\section*{Decision Trees}
\stepcounter{chapter}
\section*{Neural Networks}
\stepcounter{chapter}
\input Problems_Neural_Networks
\section*{Posterior Distributions}
\stepcounter{chapter}
\section*{Clustering}
\stepcounter{chapter}
\input Problems_Clustering
\section*{Factor Analysis}
\stepcounter{chapter}
\input Problems_Factor_Analysis
\section*{Manifold Learning}
\stepcounter{chapter}
\input Problems_Manifold_learning
\section*{Generalization Bounds}
\stepcounter{chapter}
\input Problems_Concentration


\end{document}
